\section{Related Work}

\subsection{From RLHF to RLVR}

Modern language-model post-training inherits its optimization machinery from
policy-gradient reinforcement learning, from REINFORCE \cite{williams1992reinforce} to
trust-region and clipped-surrogate methods such as TRPO and PPO
\cite{schulman2015trpo,schulman2017ppo}. Advantage estimation, especially GAE, became
the standard variance-reduction device in continuous-control RL by learning value
functions over states or prefixes \cite{schulman2015gae}. In language modeling, these
ideas entered mainstream use through RLHF systems that optimize sequence-level rewards
derived from human preferences or reward models
\cite{christiano2017preferences,stiennon2020summarize,ouyang2022instruct}.

The recent reasoning-model wave shifted part of the field from preference-based RLHF
toward reinforcement learning with verifiable rewards (RLVR), where supervision is often
a deterministic correctness signal on math or code tasks. DeepSeekMath introduced GRPO
as a practical critic-free alternative for these settings \cite{shao2024deepseekmath},
and DeepSeek-R1 popularized large-scale RL-only or RL-dominant reasoning pipelines
\cite{deepseekai2025deepseekr1}. Subsequent surveys document how quickly this regime
became the default template for open reasoning-model replication and extension
\cite{zhang2025hundreddays}. This transition matters for our setting because RLVR makes
sparse outcome rewards and long reasoning trajectories central, thereby exposing the
credit-assignment problem more directly than earlier short-form RLHF tasks.

\subsection{Critic-Based and Critic-Free Policy Optimization}

The classical solution to delayed reward is to learn a critic or value function and
convert returns into token- or prefix-level advantages. PPO-style RLHF pipelines follow
this recipe, but the value-estimation problem is particularly awkward for text
generation because prefixes are high-dimensional, non-Markov, and semantically
heterogeneous. VinePPO demonstrated that standard value heads produce poor estimates of
expected returns for reasoning tasks and proposed Monte Carlo vine-style rollout
estimates as a more accurate alternative \cite{kazemnejad2024vineppo}. GenAC extends
this idea by replacing scalar value prediction with a generative critic that uses
chain-of-thought reasoning for value estimation \cite{genac2026}.

A growing body of work questions whether value heads are necessary at all in LLM
post-training. ReMax replaces learned critics with a greedy baseline and shows that much
of PPO's practical benefit can be retained with a simpler REINFORCE-style objective
\cite{li2023remax}. Back to Basics systematically compares PPO with RLOO-style
estimators and finds that critic-free methods can match or exceed value-based baselines
in RLHF \cite{ahmadian2024backtobasics}. REINFORCE++ further strengthens this line by
improving robustness to prompt and reward-model variation without reintroducing a
learned value function \cite{hu2025reinforcepp}. In reasoning-focused RLVR, GRPO
normalizes rewards across a group of sampled responses rather than learning prefix
values \cite{shao2024deepseekmath}, and later work refines this family with theoretical
analyses, off-policy variants, and sequence-level formulations such as GSPO
\cite{rigotti2025revisiting,zheng2025gspo}. These methods substantially simplify the
optimization stack, but they generally still broadcast a trajectory-level signal across
all tokens once a scalar baseline is fixed.

\subsection{Reasoning-Specific Analyses of RLVR}

Another recent line asks why critic-free RLVR works so well for reasoning in the first
place. A Minimalist Approach to LLM Reasoning argues that strong reasoning gains can
emerge from a comparatively simple recipe that moves from rejection sampling to
REINFORCE-style online updates, suggesting that some of the field's complexity is
optional rather than essential \cite{xiong2025minimalist}. Complementarily, Wen et al.\
analyze RLVR theoretically and empirically, arguing that answer-level verifiable rewards
can nonetheless incentivize correct intermediate reasoning early in a trajectory
\cite{wen2025rlvrimplicit}. These results are important because they show that uniform
outcome-level objectives already contain nontrivial reasoning signal.

At the same time, they do not eliminate the question of \emph{which} tokens should
absorb that signal. Recent empirical analysis of RLVR training dynamics suggests that
improvements are disproportionately driven by a minority of high-entropy ``forking''
tokens, and that restricting updates to this subset can remain competitive or even
outperform full-token updates in some settings \cite{wang2025beyond8020}. Building on
this observation, EAPO adapts conditional mutual information to the autoregressive RLVR
setting and proves that the credit a token can carry is upper-bounded by its entropy,
motivating an entropy-aware modulation of per-token learning signals
\cite{wang2026eapo}. ERPO similarly identifies ``critical decision pivots'' at
high-entropy states and applies targeted exploration at those positions
\cite{chen2026erpo}. These entropy-based methods are among the closest concurrent work
to our entropy-reduction weighting: they demonstrate that entropy-guided credit
assignment can improve reasoning RL when applied with sufficient theoretical grounding.
Our work differs by systematically evaluating multiple intrinsic signals (surprisal,
entropy reduction, policy divergence) under a unified framework and documenting the
conditions under which they do and do not differentiate from uniform weighting.

\subsection{Fine-Grained Credit Assignment Beyond Uniform Token Updates}

A large parallel literature attempts to make supervision denser than a single
end-of-trajectory reward. One family learns explicit token- or process-level reward
estimators. Preference-grounded Token-level Guidance derives token-level signals from
preference data for fine-tuning \cite{yang2023pgtg}, while Discriminative Policy
Optimization learns token-level Q-style reward models from preferences without requiring
fine-grained annotation \cite{chen2025dpoqrm}. PRIME pushes this direction further for
reasoning tasks by constructing implicit process rewards and updating process reward
models online from rollouts and outcome labels \cite{cui2025prime}. Reinforced Token
Optimization (RTO) frames RLHF as a token-level MDP and derives an optimal
Q-function-based policy from the preference model, producing per-token credit without a
separate value head \cite{zhong2025rto}. These methods can produce rich token-level
feedback, but they rely on an auxiliary reward-modeling component that our approach
intentionally avoids.

Another family redistributes outcome rewards algorithmically rather than by training a
dense reward model. RED derives token-level rewards by redistributing holistic
reward-model scores over a trajectory \cite{huang2024red}. SCAR uses Shapley-inspired
attribution to estimate token or span contributions from sequence-level feedback
\cite{cao2025scar,shapley1953value}. TEMPO uses groups of sampled solutions to build a
prefix tree and compute nonparametric prefix values, yielding branch-sensitive
corrections without a learned critic \cite{tran2025tempo}. GRPO-$\lambda$ introduces
eligibility traces and lambda-returns into critic-free RLVR to propagate outcome
information backward along the sequence \cite{parthasarathi2025grpolambda}. OTB derives
an optimal token-level baseline as the ratio of expected squared advantage to expected
advantage at each position, showing theoretically and empirically that position-dependent
baselines can significantly reduce gradient variance in long-horizon LLM-RL
\cite{li2026otb}. PSPO applies potential-based reward shaping theory to construct
dense token rewards from outcome signals without altering the optimal policy
\cite{zhang2026pspo}. Compared with these approaches, our focus is narrower and cheaper:
we do not train auxiliary reward models, estimate Shapley values, construct trees, or
solve for optimal baselines, but instead ask how far one can get by reweighting the
existing policy-gradient terms using intrinsic distributional signals already present in
the forward pass.

\subsection{Alternative Action Granularities and Token Selection}

Several papers attack the same underlying problem by changing the optimization unit
rather than by redesigning the reward. MA-RLHF introduces macro actions, grouping tokens
into higher-level units to shorten the temporal distance between decisions and rewards
\cite{chai2025marlhf}. GSPO moves importance weighting and clipping from the token level
to the sequence level for greater stability in large-scale RL training
\cite{zheng2025gspo}. These methods reinforce the broader point that the granularity of
optimization matters as much as the reward definition.

Outside on-policy RL proper, token-selective optimization has also been explored in
adjacent paradigms. ConfPO emphasizes preference-critical tokens based on policy
confidence in preference optimization \cite{yoon2025confpo}, and token weighting for
long-range language modeling shows that non-uniform token emphasis can improve
optimization even in supervised settings \cite{helm2025tokenweighting}. Together with
high-entropy token analyses in RLVR \cite{wang2025beyond8020}, these results strongly
suggest that uniform token treatment is an arbitrary design choice rather than a
necessity.

\subsection{Positioning of This Work}

The literature now makes three points clear. First, critics are not strictly required
for effective LLM post-training: critic-free estimators such as RLOO, ReMax, GRPO, and
GSPO are already competitive in both RLHF and RLVR
\cite{ahmadian2024backtobasics,li2023remax,shao2024deepseekmath,zheng2025gspo}. Second,
many researchers have concluded that trajectory-level rewards are too coarse and have
introduced denser supervision through process reward models, redistribution rules, tree
structures, optimal baselines, or temporal traces
\cite{huang2024red,cui2025prime,tran2025tempo,cao2025scar,parthasarathi2025grpolambda,li2026otb,zhang2026pspo}.
Third, recent reasoning-focused analyses and entropy-aware methods indicate that only a
subset of tokens may be responsible for most of the useful RL signal, and that
entropy-guided modulation of per-token credit can improve training
\cite{wang2025beyond8020,wang2026eapo,chen2026erpo}.

Our paper sits at the intersection of these trends but occupies a distinct point in the
design space. We study \emph{intrinsic} token weighting for on-policy language RL:
credit is redistributed at the token level using signals derived from the policy's own
predictive distribution---surprisal, entropy reduction, and policy divergence---without
learned value heads, external reward models, process annotations, tree construction,
Shapley-style counterfactual attribution, or corrective rollouts. While concurrent work
such as EAPO \cite{wang2026eapo} and ERPO \cite{chen2026erpo} has independently
validated the utility of entropy as a credit signal, our contribution is complementary:
we provide a unified framework for comparing multiple intrinsic signals and a systematic
empirical analysis of when and why these signals differentiate from uniform weighting.
Rather than proposing a single method and claiming superiority, we map the space of
cheap, forward-pass-only credit signals and characterize the conditions (model
adaptation method, training regime, task difficulty) under which they do and do not
provide meaningful structure beyond a uniform baseline.
